\documentclass[10pt]{article}
\usepackage{compresr}
\cprSetHeader{Python · LangChain Integration}

\begin{document}

\cprTitleBlock{Compresr SDK · Python}{LangChain Integration}{%
    Auto-compress tool outputs in LangChain agents to cut tokens, latency,
    and cost — drop-in middleware, no app rewrites.}

% ===========================================================================
\section{Why teams use it}

Agents waste tokens on bulky tool outputs — search results, CRM dumps,
scraped pages, SQL rows — that the LLM only needs a small slice of.
Compresr rewrites each tool output \emph{against the user's query} before
it re-enters the conversation, so the model sees only what matters.

\begin{itemize}
    \item \textbf{Lower bills} — typically 40–80\% fewer tokens per tool turn.
    \item \textbf{Faster turns} — smaller prompts mean faster LLM responses.
    \item \textbf{Longer agent loops} — stay under context limits on
          multi-step runs.
    \item \textbf{Zero refactor} — one line in
          \cprInlineCode{create\_agent}; your tools stay untouched.
\end{itemize}

Compresr fails open by default: if the compression API is unreachable,
your agent keeps running on the original tool output — never a hard stop.

% ===========================================================================
\section{Install}

\begin{shcode}
pip install 'compresr[langchain]'
export COMPRESR_API_KEY=cmp_...
\end{shcode}

Requires LangChain 1.0+ and Python 3.10+.

% ===========================================================================
\section{Drop-in middleware}

Pass \cprInlineCode{CompresrToolMiddleware} to \cprInlineCode{create\_agent}.
Every tool output is compressed against the user's query before the next
LLM turn — no tool code changes.

\begin{pycode}
from langchain.agents import create_agent
from langchain_core.tools import tool
from langchain_openai import ChatOpenAI
from compresr.integrations.langchain import CompresrToolMiddleware

@tool
def web_search(query: str) -> str:
    # your real search call here
    return "<search results...>"

agent = create_agent(
    model=ChatOpenAI(model="gpt-4o-mini"),
    tools=[web_search],
    middleware=[CompresrToolMiddleware(
        target_compression_ratio=0.6,   # remove ~60% of tokens
        min_tokens=200,                 # leave short outputs alone
        query_arg="query",              # use this tool arg as the query
    )],
)

result = agent.invoke({"messages": [
    {"role": "user", "content": "Find recent pricing complaints."},
]})
\end{pycode}

That's it. Tool outputs flow through Compresr; everything else — message
shape, tool-call IDs, error handling — is preserved.

% ===========================================================================
\section{Wrap a single tool}

When you only want to compress one noisy tool (and skip the middleware
entirely), wrap it directly:

\begin{pycode}
from langchain_core.tools import tool
from compresr.integrations.langchain import wrap_tool_with_compression

@tool
def search_crm(query: str) -> str:
    """Search the CRM for customer notes."""
    return "<CRM rows...>"

search_crm = wrap_tool_with_compression(
    search_crm,
    target_compression_ratio=0.7,
    query_arg="query",
)
\end{pycode}

Works with plain LCEL, custom LangGraph graphs, or anywhere a
\cprInlineCode{StructuredTool} is used. Both sync and async paths are
covered.

% ===========================================================================
\section{Key parameters}

\renewcommand{\arraystretch}{1.15}
\begin{longtable}{@{}p{4.2cm} p{2.4cm} p{7.2cm}@{}}
\toprule
\textbf{Parameter} & \textbf{Default} & \textbf{Description} \\
\midrule
\endhead

\cprInlineCode{api\_key} & \cprInlineCode{None} &
Compresr API key. Falls back to \cprInlineCode{COMPRESR\_API\_KEY}. \\

\cprInlineCode{target\_compression\_ratio} & \cprInlineCode{0.5} &
Compression strength. $\in [0,1]$ = fraction of tokens to remove
(\cprInlineCode{0.6} = remove 60\%); $> 1$ = Nx factor
(\cprInlineCode{4} = 4x smaller). \\

\cprInlineCode{min\_tokens} & \cprInlineCode{200} &
Skip outputs below this token count — don't pay to compress what's
already small. \\

\cprInlineCode{compression\_model} & \cprInlineCode{latte\_v1} &
Compression model. \cprInlineCode{latte\_v1} is the default;
\cprInlineCode{latte\_v2} is the faster variant. \\

\cprInlineCode{query} & \cprInlineCode{None} &
Static query string — useful when every call shares the same intent. \\

\cprInlineCode{query\_arg} & \cprInlineCode{None} &
Name of the tool argument whose value should be used as the query
(e.g.\ \cprInlineCode{"query"}, \cprInlineCode{"question"}). \\

\cprInlineCode{query\_extractor} & \cprInlineCode{None} &
Callable for custom query logic — return \cprInlineCode{None} to fall
back to the default. \\

\bottomrule
\end{longtable}

\textbf{Query resolution order:} \cprInlineCode{query} →
\cprInlineCode{query\_extractor} → \cprInlineCode{query\_arg} → smart
default (most recent user message). You only need to set one.

\textit{Also available:}
\cprInlineCode{CompresrSummarizationMiddleware} (KV-cache-friendly history
summarization) and \cprInlineCode{CompresrPromptMiddleware} (last-mile
prompt budget enforcement). A
\cprInlineCode{BaseDocumentCompressor} for retriever-side compression
ships as \cprInlineCode{CompresrExtractor}.

\vfill
{\footnotesize\color{cprMuted}%
Questions or a custom deployment?
Reach the team at \href{mailto:support@compresr.ai}{support@compresr.ai}.}

\end{document}
